% Dwa skonczone stadia konstrukcji rogatej sfery Alexandera.
% Kazda para rogów wyrasta z jednego zakonczonego rogu poprzedniego
% stadium; zadna z narysowanych krzywych nie oznacza nowego 1-uchwytu.
\begin{tikzpicture}[line cap=round,line join=round,>=Latex]
  \definecolor{ahsoutline}{RGB}{38,69,83}
  \definecolor{ahsbody}{RGB}{131,178,185}
  \definecolor{ahshornA}{RGB}{103,162,174}
  \definecolor{ahshornB}{RGB}{69,124,145}
  \def\AHStube#1#2#3#4{%
    \draw[ahsoutline,line width=#2] #1;
    \draw[#4,line width=#3] #1;
  }
  \def\AHSbody{%
    \shade[inner color=ahsbody!32!white,outer color=ahsbody!84!ahsoutline]
      (0,-.57) circle (1.53);
    \draw[ahsoutline!80,line width=.7pt] (0,-.57) circle (1.53);
  }
  \def\AHSfirstpair{%
    % Koncowki dwoch wypuklosci zamieniaja strony w rzucie.
    \AHStube{(-1.18,.37)
      .. controls (-1.48,.66) and (-1.66,1.50) .. (-1.47,2.10)
      .. controls (-1.24,2.82) and (-.47,2.92) .. (.66,2.36)}
      {14pt}{11pt}{ahshornA}
    \AHStube{(1.18,.37)
      .. controls (1.48,.66) and (1.66,1.50) .. (1.47,2.10)
      .. controls (1.24,2.82) and (.47,2.92) .. (-.66,2.36)}
      {14pt}{11pt}{ahshornB}
  }
  \def\AHSchildpair{%
    % Jeden koniec rozgalezia sie na dwa mniejsze, haczykowate rogi.
    \AHStube{(0,0)
      .. controls (-.32,.19) and (-.76,.33) .. (-.79,.78)
      .. controls (-.83,1.20) and (-.29,1.48) .. (.28,1.13)}
      {6.4pt}{4.8pt}{ahshornA}
    \AHStube{(0,0)
      .. controls (.32,.19) and (.76,.33) .. (.79,.78)
      .. controls (.83,1.20) and (.29,1.48) .. (-.28,1.13)}
      {6.4pt}{4.8pt}{ahshornB}
    \filldraw[fill=ahshornA!58!white,draw=ahsoutline,line width=.45pt]
      (.28,1.13) circle (.095);
    \filldraw[fill=ahshornB!58!white,draw=ahsoutline,line width=.45pt]
      (-.28,1.13) circle (.095);
  }
  \begin{scope}[shift={(-2.85,0)},scale=.90]
    \AHSbody
    \AHSfirstpair
    \filldraw[fill=ahshornA!58!white,draw=ahsoutline,line width=.55pt]
      (.66,2.36) circle (.085);
    \filldraw[fill=ahshornB!58!white,draw=ahsoutline,line width=.55pt]
      (-.66,2.36) circle (.085);
  \end{scope}
  \node[font=\small] at (-2.85,-2.33) {$n=1$: dwa rogi};
  \draw[->,ahsoutline,line width=.8pt] (-.81,.75) -- (.81,.75);
  \node[font=\scriptsize,align=center,text=ahsoutline]
    at (0,1.13) {rozgałęzienie\\każdej końcówki};
  \begin{scope}[shift={(2.85,0)},scale=.90]
    \AHSbody
    \AHSfirstpair
    \begin{scope}[shift={(-.66,2.36)},rotate=12,scale=.55]
      \AHSchildpair
    \end{scope}
    \begin{scope}[shift={(.66,2.36)},rotate=-12,scale=.55]
      \AHSchildpair
    \end{scope}
  \end{scope}
  \node[font=\small] at (2.85,-2.33) {$n=2$: cztery końcówki};
\end{tikzpicture}
